% Generated draft: requires mathematical and source review.
\documentclass[11pt]{article}
\usepackage[margin=25mm]{geometry}
\usepackage{fontspec}
\setmainfont{Latin Modern Roman}
\usepackage{amsmath,amssymb,mathtools,unicode-math}
\setmathfont{Latin Modern Math}
\usepackage{enumitem,needspace,xurl,hyperref}
\setlength{\parindent}{0pt}
\setlength{\parskip}{5pt}
\setlength{\emergencystretch}{3em}
\newcommand{\N}{\mathbb{N}}
\newcommand{\Z}{\mathbb{Z}}
\newcommand{\Q}{\mathbb{Q}}
\newcommand{\R}{\mathbb{R}}
\newcommand{\C}{\mathbb{C}}
\newcommand{\F}{\mathbb{F}}
\newcommand{\A}{\mathbb{A}}
\newcommand{\PP}{\mathbb{P}}
\newcommand{\E}{\mathbb{E}}
\newcommand{\eps}{\varepsilon}
\newcommand{\dd}{\,\mathrm{d}}
\begin{document}
\section*{UnsolvedMath ID: 30004075}
\textbf{Source problem code:} catalogue entry
\subsection{Definitions and mathematical statement}
Let \(p\) be a prime and let \(R\) be a commutative Noetherian integral domain of characteristic \(p\). The ring \(R\) is \(F\)-finite if its Frobenius map \(F:R\to R\), \(r\mapsto r^p\), is finite. Let \(0\ne\mathfrak a\subseteq R\) be an ideal and let \(t\in\R_{>0}\).

We use the following convention for the finitistic test ideal of the pair \((R,\mathfrak a^t)\). For \(e\geq0\), put \(q=p^e\), and let \(R^{(e)}\) denote the bimodule with usual left multiplication and right multiplication \(u\cdot r=ur^q\). For an \(R\)-module \(M\), put
\[
F_R^e(M)=R^{(e)}\otimes_R M,\qquad z^{[q]}=1\otimes z.
\]
For each finitely generated \(R\)-module \(M\), define
\[
0_M^{*\mathfrak a^t}
=
\left\{z\in M:\ \exists c\in R\setminus\{0\}\ \exists e_0\geq0\ \forall e\geq e_0,\quad
c\,\mathfrak a^{\lceil t p^e\rceil}z^{[p^e]}=0
\text{ in }F_R^e(M)\right\}.
\]
Here the displayed vanishing means that every element of \(c\mathfrak a^{\lceil tp^e\rceil}\) annihilates \(z^{[p^e]}\). Define
\[
\tau_{\mathrm{fg},R}(\mathfrak a^t)
=
\bigcap_{M\ \text{finitely generated over }R}
\operatorname{Ann}_R\!\left(0_M^{*\mathfrak a^t}\right).
\]
The notation \(\mathfrak a^t\) denotes the ideal together with its real exponent in this definition; it does not assert that an ordinary ideal power with nonintegral exponent has been defined.

\textbf{Localization question.} For every such \(F\)-finite ring \(R\), nonzero ideal \(\mathfrak a\), positive real number \(t\), and multiplicatively closed subset \(S\subseteq R\) with \(1\in S\) and \(0\notin S\), does the equality of ideals in \(S^{-1}R\)
\[
S^{-1}\tau_{\mathrm{fg},R}(\mathfrak a^t)
=
\tau_{\mathrm{fg},S^{-1}R}\!\left((S^{-1}\mathfrak a)^t\right)
\]
hold? The right-hand ideal is defined by the same construction over \(S^{-1}R\), using all finitely generated \(S^{-1}R\)-modules.

\textbf{Unresolved source qualifications.} The supplied record does not explicitly define \(\tau_{\mathrm{fg}}\), specify the ambient domain of \(t\), or state whether \(0\) is excluded from \(S\). The formulation above explicitly uses the finitistic generalized tight-closure definition, positive real exponents, and nonzero localization. These conventions require confirmation against the cited source. If \(0\in S\) is intended, the localization is the zero ring, outside the domain convention used above, and the supplied record gives no corresponding test-ideal convention. The record supplies no additional hypotheses or precise special cases for its qualification that the question remains unresolved beyond additional hypotheses.
\subsection{Short English statement}
Does the finitistic test ideal of every pair \((R,\mathfrak a^t)\), with \(R\) an \(F\)-finite Noetherian domain of positive characteristic, commute with localization? The displayed formulation makes explicit the conventions for generalized tight closure, positive real exponents, and multiplicative sets avoiding zero. The supplied open-status assessment is dated curation requiring independent review.
\subsection{Sources}
\begin{itemize}
\item Local registry: \texttt{data/unsolvedmath/problems.json}, supplied record with permanent UnsolvedMath ID \(30004075\), title ``Localization of Finitistic Test Ideals,'' exact source problem code \texttt{catalogue entry}. The record belongs to the Oberwolfach Reports Open Problems collection and supplies proposed year 2019; no individual proposer is identified.
\item Supplied report citation: \emph{Singularities and Homological Aspects of Commutative Algebra} (2020), Oberwolfach Reports, pp.~407--464, DOI \texttt{10.4171/owr/2019/6}. \url{https://doi.org/10.4171/owr/2019/6}. The literature-source list also labels this reference ``Oberwolfach Report 6/2019.'' The supplied 2020 citation year and 2019 report designation are retained without reconciliation.
\item Supplied literature source, labeled ``St\"uckrad--Vogel resolution'': \url{https://doi.org/10.1016/j.aim.2019.02.029}. The record's assessment says that the St\"uckrad--Vogel conjecture mentioned in its provenance is solved; it distinguishes this from the extracted localization question.
\item Supplied literature source, labeled ``Lyubeznik--Smith on localization'': \url{https://doi.org/10.1090/S0002-9947-01-02643-5}. The supplied record does not identify the exact results or hypotheses relevant to the displayed equality.
\item Literature assessment checked at \texttt{2026-08-22}: the supplied status is ``open,'' with the qualification that the extracted localization compatibility question for finitistic test ideals remains unresolved beyond additional hypotheses. This is supplied dated curation derived from the OWR record, not an independently verified assessment of the literature. The source explicitly requests independent verification of time-sensitive status and cited claims.
\item The registry supplies the statement-status label \texttt{corrected\_verified} and describes notation, typography, duplicate-clause, or definition repairs as preserving the mathematical content. These are supplied editorial assertions, not independent mathematical verification of this draft or its conventions.
\end{itemize}
\end{document}
